\documentclass[11pt,a4paper]{article}
\usepackage[utf8]{inputenc}
\usepackage[T1]{fontenc}
\usepackage{lmodern}
\usepackage{amsmath,amssymb}
\usepackage{booktabs}
\usepackage{geometry}
\geometry{margin=2.4cm}
\usepackage[hidelinks]{hyperref}
\usepackage[numbers,sort&compress]{natbib}
\usepackage{enumitem}
\setlist{nosep}

\newcommand{\verif}[1]{\textsf{\small[#1]}}
\newcommand{\takes}{\par\noindent\textbf{What we take.}\ }
\newcommand{\used}{\par\noindent\textbf{Where it is used.}\ }

\title{Credible procedural weather without simulation:\\
a literature review for WorldBuilder's synoptic weather model}
\author{WorldBuilder --- branch \texttt{weather-synoptic}}
\date{19 November 2026}

\begin{document}
\maketitle

\begin{abstract}
WorldBuilder renders volumetric clouds over procedural planets whose radii range from
600\,km to Jupiter scale. Its weather must look credible from orbit, from an aircraft and
from the ground, cost nothing per frame, and accept local presets pinned at coordinates.
This review collects the meteorological and graphics literature behind the chosen design:
an \emph{analytic synoptic-object model}, in which jets are Rossby waves, lows, highs and
tropical cyclones are parametric objects on closed-form tracks, and pressure, wind, fronts
and cloud genera are derived from them by conceptual-model rules. For each source we state
what it says, the quantitative relations we take from it, and where they enter the design.
Verification tags give the level at which each claim was checked:
\verif{FT} full text, \verif{abs} abstract, \verif{2nd} secondary source only,
\verif{meta} bibliographic data only.
\end{abstract}

\tableofcontents

%====================================================================
\section{Requirements that drive the review}
\begin{enumerate}
\item \textbf{Recognisable structure}: jets that meander; cyclones with fronts and the
  classic comma cloud; anticyclones that are clear or capped by stratocumulus; tropical
  cyclones with an eye.
\item \textbf{No simulation}: the weather state $W(t)$ must be a pure function of time and
  seed. For now it is evaluated at one fixed $t$.
\item \textbf{Scales from the planet}: cyclone size, number of jets and belt latitudes
  follow the radius, the rotation and the stratification. They are not resized Earth tables.
\item \textbf{Local presets}: a situation pinned at a coordinate (``warm front
  approaching'') must sit naturally inside the global pattern.
\end{enumerate}

%====================================================================
\section{General circulation and its scaling with rotation}

\paragraph{Held \& Hou (1980) \cite{held1980}.}
An axisymmetric, angular-momentum-conserving theory of the Hadley cell. It predicts the
cell's width, its heat flux, the latitude of the subtropical jet and the surface
easterly/westerly pattern \verif{abs}. In the small-angle limit the edge is
\begin{equation}
  \phi_H = \left(\frac{5}{3}\,\frac{gH\,\Delta\theta}{\Omega^2 a^2\,\theta_0}\right)^{1/2}
  \;\propto\; \Omega^{-1},
\end{equation}
as quoted in \citet[Eq.~56]{showman2009} \verif{FT of the secondary}.
\takes The functional dependence on $\Omega a$ and the thermal Rossby number
$Ro_T \propto gH\Delta\theta/(\Omega a)^2$. Its exponent $1/2$ is used only as the upper
bound.
\used \texttt{Scales.hadley\_edge}.

\paragraph{Schneider (2006) \cite{schneider2006}.}
A review of the general circulation. It sets the Held--Hou bound
$\phi_M \propto (gH\Delta_h/\Omega^2a^2T_0)^{1/2}$ against the edge limited by baroclinic
instability, $\phi_c \approx (gH\Delta_v/\Omega^2a^2T_0)^{1/4} \propto \Omega^{-1/2}$
(Eq.~3, after Held 2000) \verif{FT}. Earth's observed cell is closer to the eddy-limited
regime.
\takes $\phi_H = 30^\circ\,(Ro_T/Ro_{T,\oplus})^{1/4}$, calibrated to Earth's
$30^\circ$ and clamped to $[10^\circ, 60^\circ]$.
\used \texttt{Scales.hadley\_edge}, and through it every belt latitude.

\paragraph{Walker \& Schneider (2006) \cite{walker2006}.}
An idealised dry GCM varied over radius, rotation and forcing. The Hadley strength and
extent follow regimes that axisymmetric theory misses, with upper-branch Rossby numbers of
0.1--0.8 \verif{abs}. Their extent falls with $\Omega$, but more slowly than $\Omega^{-1}$
(per \cite{schneider2006}, \verif{FT}).
\takes This justifies the $\Omega^{-1/2}$ choice above instead of Held--Hou's $\Omega^{-1}$.

\paragraph{Kaspi \& Showman (2015) \cite{kaspi2015}.}
An idealised moist GCM run from $1/16$ to $8\,\Omega_\oplus$, with radius, gravity, mass
and stellar flux also varied. At slow rotation the Hadley cell is nearly global, and it
contracts as rotation speeds up. Jets multiply with rotation, and fast rotators have larger
equator-to-pole temperature differences \verif{FT}.
\takes These are qualitative checks on our scalings. A slowly rotating small world (our
600\,km planet with a 6\,h day has $\Omega a/NH \approx 1.75$ against Earth's $4.65$)
should show a broad Hadley cell and a single jet. A fast, large world should show many jets.
\used Test expectations in \texttt{tests/test\_weather\_scales.py}.

\paragraph{Showman, Cho \& Menou (2009) \cite{showman2009}.}
A review of atmospheric circulation on exoplanets. It notes that the number of jets scales
as $\Omega^{1/2}$ and the Rhines scale as $\approx \pi (U/\beta)^{1/2}$ \verif{FT}.
\takes This cross-checks the jet count.

%====================================================================
\section{Jets and Rossby waves}

\paragraph{Rossby et al.\ (1939) \cite{rossby1939}; Rossby (1945) \cite{rossby1945}.}
The 1939 paper gives the barotropic phase speed
\begin{equation}
  c = U - \frac{\beta}{K^2}, \qquad L_s = 2\pi\sqrt{U/\beta},
\end{equation}
where $L_s$ is the stationary wavelength \verif{2nd; the 1939 PDF was not reachable}. The
1945 paper adds group velocity and energy propagation \verif{meta}.
\takes Each jet axis is
$\phi_j(\lambda,t)=\phi_0+\sum_m A_m\sin(k_m\lambda-\omega_m t+\theta_m)$ with
$\omega_m = k_m(U - \beta/K_m^2)/(a\cos\phi_0)$. Long waves come out quasi-stationary and
short ones move east, which is the behaviour a forecaster expects to see.
\used \texttt{synoptic/jets.py}.

\paragraph{Rhines (1975) \cite{rhines1975}.}
On a $\beta$-plane, the inverse energy cascade stops at $k_\beta=(\beta/2U)^{1/2}$, with $U$
the rms particle speed. Beyond it the flow turns into waves and then into alternating zonal
jets. $k_\beta^{-1}$ is about 1000\,km in the atmosphere \verif{abs}.

\paragraph{Williams (1978) \cite{williams1978}.}
A forced barotropic model on the sphere reproduces Jovian banding with a band scale of
$L_\beta=\pi(2U/\beta)^{1/2}$. With Earth's parameters it gives one midlatitude westerly jet
and the trades \verif{abs}.
\takes Jet spacing $L_\beta=\pi\sqrt{2U/\beta_{45}}$ with $U\approx20$\,m\,s$^{-1}$. The
number of eddy-driven jets is the width of the extratropics divided by $L_\beta$, clamped to
1--8.
\used \texttt{Scales.jet\_lats}.

\paragraph{Uccellini \& Kocin (1987) \cite{uccellini1987}.}
In eight East-Coast heavy-snow cases, a direct circulation in the entrance of one jet
streak couples with the indirect circulation in the exit of another. This sets up
frontogenesis and cyclogenesis \verif{abs}.
\takes The placement rule: new lows spawn on the downstream side of a trough, under the
jet's left exit. Pinned lows get jet phases refitted so that this still holds.
\used \texttt{systems.population}, \texttt{presets.solve}.

\paragraph{Hoskins \& Valdes (1990) \cite{hoskins1990}; Hoskins \& Hodges (2002) \cite{hoskins2002}.}
Storm tracks maintain themselves: diabatic heating restores the baroclinicity that eddy fluxes
remove \verif{abs}. Feature tracking on reanalyses shows where the tracks lie, and shows that
anticyclonic features are weaker and less coherent than cyclones \verif{abs}.
\takes Storm tracks are bands just poleward of the Hadley edge, starting downstream of
continents. Highs get broader, weaker kernels than lows.
\used \texttt{Scales.storm\_track}, \texttt{fields.pressure}.

\paragraph{Rex (1950a,b) \cite{rex1950a,rex1950b}.}
The definition of blocking: the westerlies split into two branches over at least
$45^\circ$ of longitude, persisting at least 10 days \verif{2nd}.
\takes About 10\,\% of migratory highs are blocking highs: three times the lifetime, a fifth
of the speed, and they push the local jet poleward.
\used \texttt{systems.py}.

%====================================================================
\section{Extratropical cyclones and fronts}

\paragraph{Bjerknes \& Solberg (1922) \cite{bjerknes1922}.}
The Norwegian model: a wave on the polar front grows into a cyclone with warm and cold
fronts, which then occludes \verif{meta, through a 1923 review}.
\takes The life-cycle stages: incipient wave, open wave, mature, occluding,
occluded or decaying. Each is mapped to the stage parameter $s\in[0,1]$.
\used \texttt{systems.stage}, \texttt{fronts.fronts}.

\paragraph{Shapiro \& Keyser (1990) \cite{shapiro1990}.}
An alternative life cycle typical of intense marine cyclones: frontal fracture, a bent-back
warm front with a ``T-bone'', then warm-core seclusion \verif{2nd}.
\takes A seeded share of lows (\texttt{variant="sk"}) follows this geometry instead.
\used \texttt{fronts.py}.

\paragraph{Schultz \& Vaughan (2011) \cite{schultz2011}.}
Occlusion comes from the warm tongue wrapping up by deformation and rotation, not from the
cold front catching up. Warm-type occlusions are more common. Cyclones often keep deepening
by 10--30\,hPa for 12--36\,h after the occlusion forms \verif{abs}.
\takes The occluded front wraps poleward around the centre as $s$ passes
$\approx0.5$. Peak intensity comes after the occlusion starts, so the strength curve peaks
at $s\approx0.6$.
\used \texttt{fronts.py}, \texttt{systems.py}.

\paragraph{Carlson (1980) \cite{carlson1980}.}
An isentropic relative-flow analysis. The warm conveyor belt, the cold conveyor belt and the
dry airstream organise the comma cloud, and the boundaries between airstreams give its sharp
edges \verif{abs}.
\takes The zones: the warm conveyor belt ascends over the warm front (the frontal cloud
shield), the cold conveyor belt feeds the comma head, and the dry intrusion makes a clear
dry slot behind the cold front. Zone edges are sharp (narrow smoothstep), not Gaussian.
\used \texttt{rules.py} zone table.

\paragraph{Browning (1990) \cite{browning1990}; Browning \& Roberts (1994) \cite{browning1994}.}
A review of how cloud and precipitation are organised in extratropical cyclones
\verif{meta}. The 1994 case study shows the cloud head's sharp convex edge, the warm conveyor
belt splitting under the dry intrusion, and line convection on the cold front \verif{abs}.
\takes Narrow line convection, Cb with congestus, along the cold front. Embedded
convection in the comma head.
\used \texttt{rules.py}.

\paragraph{Warm-front cloud sequence.}
Teaching sources give the order Ci $\to$ Cs $\to$ As $\to$ Ns (then St/fog) as the surface
front approaches. The frontal slope is about 1:100--1:300, with cirrus around 1000\,km or
more ahead of the surface front on Earth \verif{2nd; sources disagree on the slope}. We cite
the textbook treatment \cite{holton2013} and the WMO atlas \cite{wmo2017} for the genera.
\takes A band of about $1\,L_R$ on the cold side of the warm front, with the genus chosen by
the distance to the surface front. Test: \texttt{test\_warm\_front\_sequence}.

%====================================================================
\section{Anticyclones and the stratocumulus decks}

\paragraph{Wood (2012) \cite{wood2012}.}
Stratocumulus needs strong lower-tropospheric stability, subsidence and a shallow coupled
boundary layer. It is driven by cloud-top longwave cooling and covers more than 50\,\% of
the cool subtropical and midlatitude oceans every year \verif{abs}.
\paragraph{Klein \& Hartmann (1993) \cite{klein1993}.}
They identify ten regions of persistent stratus, the subtropical marine ones on the eastern
sides of ocean basins. Stratus amount rises by about 6\,\% per 1\,K of lower-tropospheric
stability \verif{abs}.
\takes High cores are clear (subsidence). Over cool sea on the eastern basin flank, an Sc
deck sits under the inversion, with St/fog near the coast. Deck coverage rises with a
stability proxy (cold sea under a strong high).
\used \texttt{rules.py}; the subtropical high placement in \texttt{systems.py}
(longitude bias toward eastern basins).

%====================================================================
\section{Tropical cyclones}

\paragraph{Holland (1980) \cite{holland1980}.}
A two-parameter radial pressure profile
\begin{equation}
  p(r) = p_c + \Delta p\,\exp\!\left[-\left(R_m/r\right)^B\right],
\end{equation}
with gradient wind derived from it \verif{abs; the B range is from secondary sources}.
\paragraph{Willoughby \& Rahn (2004) \cite{willoughby2004}.}
Fits to 493 aircraft sorties give $B$ between 0.5 and 2.5, with mean
$1.31\pm0.36$ \verif{FT}.
\takes The TC pressure kernel, with $B\sim\mathcal{N}(1.3,0.35)$ clamped to $[0.8,2.2]$,
and gradient wind near the core.
\used \texttt{fields.pressure}, \texttt{fields.wind}.

\paragraph{Gray (1968) \cite{gray1968}; Camargo et al.\ (2007) \cite{camargo2007}.}
Most genesis is equatorward of $20^\circ$, on the poleward side of the equatorial trough,
where vertical shear is weak \verif{abs}. The genesis potential index
$\mathrm{GP}=|10^5\eta|^{3/2}(H/50)^3(V_{pot}/70)^3(1+0.1V_{shear})^{-2}$ combines
absolute vorticity, humidity, potential intensity and shear \verif{FT}. The
$26^\circ$C sea-temperature threshold comes from Gray's later papers, not 1968.
\takes Genesis only over sea, inside $[5^\circ, \min(25^\circ, 0.8\phi_H)]$, weighted by a
GP-like proxy built from the planet's temperature and humidity fields.
\used \texttt{Scales.tc\_band}, \texttt{systems.population}, and the HURRICANE pin check.

\paragraph{Emanuel et al.\ (2006) \cite{emanuel2006}; Bloemendaal et al.\ (2020) \cite{bloemendaal2020}.}
Synthetic tracks generated without a full model: Markov chains, or beta-and-advection
steering by synthetic winds \verif{abs}. STORM resamples 38 years of observations into
10\,000 synthetic years that preserve the statistics \verif{abs}.
\takes This is precedent that statistically generated tracks are credible. Our TCs move west
and poleward, recurve at the subtropical ridge, and become extratropical: closed-form
parabolas in the local frame.
\used \texttt{systems.position}.

%====================================================================
\section{Wind from pressure}
Standard textbook material \cite{holton2013,vallis2017} \verif{meta}:
\begin{itemize}
\item the geostrophic wind $\mathbf{v}_g = \hat{\mathbf{k}}\times\nabla p/(\rho f)$;
\item the gradient-wind correction in curved flow;
\item boundary-layer friction turns the surface wind across the isobars toward low pressure
  and slows it.
\end{itemize}
\takes The surface wind is geostrophic, turned $25^\circ$ and scaled by 0.7, and blends to
the trades for $|\phi|<10^\circ$, where $f\to0$. Its circulation is counter-clockwise around
northern lows (test: \texttt{test\_circulation\_sense}).
\used \texttt{fields.wind}. It replaces the latitude-only wind in the orographic term and the
fixed anvil axis.

%====================================================================
\section{Games and graphics practice}

\paragraph{Schneider \& Vos (2015, 2017); Schneider (2022, 2023) \cite{schneider2015,schneider2017,schneider2022,schneider2023}.}
Horizon Zero Dawn's volumetric clouds, and later Nubis, are driven by a 2D \emph{weather
map} with coverage, type and precipitation channels, authored and then animated at regional
scale. Later versions add flyable cloud environments and voxel SDF clouds
\verif{2nd and course abstracts}.
\takes This confirms the split WorldBuilder already has: a weather layer writes maps, the
renderer reads them. Our work changes only what fills the maps.

\paragraph{Hillaire (2016) \cite{hillaire2016}.}
Frostbite's physically based sky, atmosphere and clouds under dynamic weather
\verif{2nd}. This is renderer context only.

\paragraph{Bauer (2019) \cite{bauer2019}.}
Red Dead Redemption 2's integrated atmosphere and clouds. Weather is authored as regional
presets with transitions \verif{meta; content not verified}.
\takes This is the closest precedent for local presets. Its weakness, that regions are not
globally consistent, is what our object-relative presets address.

\paragraph{Webanck et al.\ (2018) \cite{webanck2018}.}
Procedural cloudscapes built from primitive-based field functions, with cloud types by
altitude, animated by optimal-transport morphing along terrain- and wind-aware paths
\verif{abs}.
\takes Genus assignment by altitude band and situation.

\paragraph{Microsoft Flight Simulator (meteoblue) \cite{meteoblue_msfs}.}
Live weather from a global forecast model (about 250 million grid boxes, 60 levels) fed to
the renderer \verif{vendor pages}.
\takes This is realism through data. It is not applicable to procedural planets, but it is
the quality bar.

%====================================================================
\section{Approaches considered and rejected}

\paragraph{Fluid simulation \cite{harris2003,goswami2017}.}
GPU cloud dynamics \verif{meta} and Lagrangian parcels for landscape-scale convection
\verif{2nd}. These work at the scale of a single cloud or a landscape. They are too
expensive and at the wrong scale for planetary synoptic weather, which is the user's
``no simulation'' constraint.

\paragraph{Pre-baked climate models: ExoPlaSim \cite{paradise2022}, Isca \cite{vallis2018}.}
Laptop-runnable GCMs for arbitrary planets, with rotation, radius and continents
configurable \verif{abs}. They are rejected for $W(t)$ for three reasons:
\begin{itemize}
\item the data volume;
\item a T21--T42 grid is a handful of cells on a 600\,km world;
\item local pins would need nudging and a rerun.
\end{itemize}
They remain a possible offline cross-check of our climatology.

\paragraph{Reanalysis analogues: ERA5 \cite{hersbach2020}.}
Hourly 31\,km global fields since 1950 \verif{abs}. Earth cut-outs pasted onto other planets
would carry Earth's scales and leave seams. ERA5 could serve as a reference for kernel shapes.

\paragraph{Stochastic weather generators: WGEN \cite{richardson1981}.}
Markov-chain wet/dry days with conditioned temperature and radiation \verif{abs}. This is a
point process with no spatial structure, so no fronts or systems.

%====================================================================
\section{Resulting scales for WorldBuilder's planets}
The relations above, with $NH=100$\,m\,s$^{-1}$, $U=20$\,m\,s$^{-1}$, Eady wavelength
$3.9\,L_R$ and $\phi_H\propto Ro_T^{1/4}$ calibrated to Earth, give the values below. They
were computed while writing this review. The values that ship are those of
\texttt{Scales.from\_planet}; the full set of relations is in \S\ref{sec:equations}.

\begin{center}
\begin{tabular}{lrrr}
\toprule
 & Kerbin-like & Earth & rocky Jupiter \\
\midrule
radius (km) & 600 & 6\,371 & 69\,911 \\
day (h) & 6 & 23.934 & 9.925 \\
$\Omega a / NH$ & 1.75 & 4.65 & 123 \\
$L_R$ at $45^\circ$ (km) & 243 & 970 & 402 \\
storm-track wavenumber $k^*$ & 2.8 & 7.5 & $\approx 200$ \\
$L_\beta$ (degrees of latitude) & 72 & 44 & 8.6 \\
Hadley edge & $49^\circ$ & $30^\circ$ & $10^\circ$ (clamped from $5.8^\circ$) \\
eddy-driven jets / hemisphere & 1 (clamped from 0.4) & 1 & $\approx 7$ \\
\bottomrule
\end{tabular}
\end{center}

The 600\,km world behaves like a slowly rotating Earth: a broad Hadley cell and two or three
large lows per hemisphere. Earth-like zonation at that radius needs $\Omega a$ equal to
Earth's, which is a day of about 2.3\,h. \texttt{day\_length\_h} and
\texttt{stratification} are the design knobs for this.

%====================================================================
\section{The model's equations}
\label{sec:equations}
This section lists every relation the implementation uses, so that the code can be checked
against it line by line. The notation is:
\begin{itemize}
\item planet radius $a$ (m), latitude $\phi$, longitude $\lambda$, time $t$ (s);
\item hemisphere sign $h=\operatorname{sign}\phi$ (north $+1$);
\item the smoothstep $S(e_0,e_1,x)=u^2(3-2u)$ with $u=\operatorname{clamp}\!\big((x-e_0)/(e_1-e_0),0,1\big)$;
\item the bump $\Pi(x;x_0,x_1,e)=S(x_0-e,x_0,x)\,\big(1-S(x_1,x_1+e,x)\big)$;
\item the Gaussian $G(x;x_0,\sigma)=\exp[-((x-x_0)/\sigma)^2]$.
\end{itemize}
Constants marked $^\dagger$ are starting values, to be tuned on preview maps within the
ranges of the cited sources. Every other relation is fixed by the design.

%--------------------------------------------------------------------
\subsection{Planetary scales}
\label{sec:eq-scales}
Inputs: the radius $a$, the day length $d$ in hours, and the stratification
$NH$ ($100$\,m\,s$^{-1}$ by default, Earth's value).
\begin{align}
  \Omega &= \frac{2\pi}{3600\,d}, &
  f(\phi) &= 2\Omega\sin\phi, &
  \beta(\phi) &= \frac{2\Omega\cos\phi}{a}, \\
  L_R &= \frac{NH}{f(45^\circ)}, &
  \lambda_E &= 3.9\,L_R, &
  k^* &= \frac{2\pi a\cos 45^\circ}{\lambda_E}.
\end{align}
$L_R$ is the deformation radius and $\lambda_E = 2\pi L_R/1.61$ is the most unstable Eady
wavelength \cite{holton2013,vallis2017}. $k^*$ is the storm-track zonal wavenumber at
$45^\circ$.

The radius of an extratropical system is
\begin{equation}
  r_c = 0.25\,\lambda_E\,\xi_r,\qquad \xi_r\sim\mathcal{U}(0.75,1.3).
\end{equation}

The jet spacing \cite{williams1978,rhines1975}, with eddy velocity $U_e=20$\,m\,s$^{-1}$, is
\begin{equation}
  L_\beta = \pi\sqrt{\frac{2U_e}{\beta(45^\circ)}},\qquad
  L_\beta^\circ = \frac{180}{\pi}\frac{L_\beta}{a}.
\end{equation}

The Hadley edge, eddy-limited, $\propto\Omega^{-1/2}$ \cite{schneider2006}, is
\begin{equation}
  Ro_T = \left(\frac{NH}{\Omega a}\right)^2,\qquad
  \phi_H = \operatorname{clamp}\!\left(30^\circ\left(\frac{Ro_T}{Ro_{T,\oplus}}\right)^{1/4},
  \,10^\circ,\,60^\circ\right),
\end{equation}
with $Ro_{T,\oplus}$ evaluated for $a=6.371\times10^6$\,m, $d=23.934$\,h and $NH=100$.

The belts, in degrees of $|\phi|$, are:
\begin{align}
  \text{storm track}\quad & [\phi_H+5^\circ,\ \min(\phi_H+30^\circ, 80^\circ)],\\
  \text{eddy-driven jets}\quad & n_j=\operatorname{clamp}\!\left(\operatorname{round}
     \frac{75^\circ-\phi_H}{L_\beta^\circ},1,8\right),\\
  & \phi_{j,i}=\phi_H+\Big(i+\tfrac12\Big)\frac{75^\circ-\phi_H}{n_j},\ i=0..n_j-1,\\
  \text{subtropical jet}\quad & \phi_{j,\mathrm{st}}=\phi_H,\\
  \text{tropical cyclone band}\quad & [5^\circ,\ \min(25^\circ,0.8\,\phi_H)].
\end{align}

The tropical-cyclone radius of maximum wind, with $V_{TC}=20$\,m\,s$^{-1}$$^\dagger$, is
\begin{equation}
  R_m = 0.1\,r_{TC},\qquad
  r_{TC}=\operatorname{clamp}\!\left(\frac{V_{TC}}{f(20^\circ)},\,0.1\,r_c,\,0.5\,r_c\right).
\end{equation}

Jet and steering speeds$^\dagger$: $U_{\mathrm{st}}=30$, $U_j=25$, and
$c_s=0.6\,U_j$ (m\,s$^{-1}$).

%--------------------------------------------------------------------
\subsection{Jet axes as Rossby waves}
\label{sec:eq-jets}
For each hemisphere and each jet $j$, there are three waves $m=1,2,3$ with integer
wavenumbers (needed for continuity in $\lambda$):
\begin{equation}
  k_m\in\Big\{\max\!\big(1,\operatorname{round}\tfrac{k^*}{2}\big),\
  \max\!\big(1,\operatorname{round}k^*\big),\
  \max\!\big(1,\operatorname{round}\tfrac{3k^*}{2}\big)\Big\}.
\end{equation}
The axis is
\begin{equation}
  \phi_j(\lambda,t)=\phi_{0,j}+h\sum_{m=1}^{3}A_m\sin\!\big(k_m\lambda-\varpi_m t+\theta_m\big),
\end{equation}
with the dispersion relation \cite{rossby1939}
\begin{equation}
  K_m=\frac{k_m}{a\cos\phi_{0,j}},\qquad
  c_m = U_j-\frac{\beta(\phi_{0,j})}{K_m^2},\qquad
  \varpi_m = \frac{k_m\,c_m}{a\cos\phi_{0,j}}.
\end{equation}
Here $\varpi_m t$ is reduced modulo $2\pi$ in float64 before evaluation. The amplitudes are
\begin{equation}
  A_m = A_{\mathrm{tot}}\,\frac{w_m}{\sum w},\quad w_m\sim\mathcal{U}(0.5,1),\qquad
  A_{\mathrm{tot}}=\min(8^\circ,\,0.25\,L_\beta^\circ),
\end{equation}
and the phases $\theta_m\sim\mathcal{U}(0,2\pi)$ are seeded, or refitted by pins
(\S\ref{sec:eq-pins}).

Troughs are the minima of $h\,\phi_j$ in $\lambda$ (equatorward excursions), and ridges are the
maxima. Both are found on a 512-point grid in $\lambda$, refined by a parabolic fit.

The upper wind has a Gaussian jet profile across the meandering axis, with half-width
$w_j=0.25\,L_\beta$$^\dagger$:
\begin{equation}
  u_\parallel = U_j\,G(y;0,w_j),\qquad
  y = a\,(\phi-\phi_j(\lambda,t)).
\end{equation}
Its stream function, with contours along the axis, is
\begin{equation}
  \psi_j = -U_j\,w_j\,\frac{\sqrt{\pi}}{2}\,\operatorname{erf}\!\left(\frac{y}{w_j}\right).
\end{equation}
The total is
$\psi=\sum_j\psi_j+\sum_i a_i\exp(-|\mathbf{d}-\mathbf{c}_i|^2\kappa_i)$, where the system
bumps keep the sign rules of today's \texttt{stream\_vortices}, and the wind is
$\mathbf{v}=\hat{\mathbf{z}}\times\nabla\psi$.

%--------------------------------------------------------------------
\subsection{Synoptic objects: tracks, stages and strength}
\label{sec:eq-objects}
\paragraph{Stage and strength.}
An object has birth time $t_b$, lifetime $T$ and peak strength $P$. Then
\begin{equation}
  s(t)=\frac{t-t_b}{T}\in[0,1],\qquad
  \Sigma(s)=P\;S(0,s_p,s)\,\big(1-S(s_p,1,s)\big),\quad s_p=0.6,
\end{equation}
where $s_p=0.6$ means peak intensity comes after the occlusion starts \cite{schultz2011}.
The lifetimes are
\begin{equation}
\begin{gathered}
  T_{\mathrm{LOW}}=2.5\,\frac{\lambda_E}{c_s}\,\xi_T,\quad
  T_{\mathrm{TC}}=1.5\,T_{\mathrm{LOW}}\,\xi_T,\quad
  T_{\mathrm{HIGH}}=1.5\,T_{\mathrm{LOW}}\,\xi_T,\\
  T_{\mathrm{block}}=3\,T_{\mathrm{HIGH}},\qquad \xi_T\sim\mathcal{U}(0.8,1.2).
\end{gathered}
\end{equation}
For Earth this gives $T_{\mathrm{LOW}}\approx 5.8$ days.

\paragraph{Tracks.}
Let $\tau=t-t_b$.
\begin{itemize}
\item LOW and CUTOFF:
\begin{equation}
  \lambda(t)=\lambda_0+\frac{\gamma\,c_s\,\tau}{a\cos\phi_0},\qquad
  \phi(t)=\phi_0+h\,\frac{v_p\,\tau}{a},\qquad v_p=0.1\,c_s,
\end{equation}
with $\gamma=1$ for a LOW and $\gamma=0.2$ for a CUTOFF. For migratory HIGHs
$\gamma=0.5$, $v_p=0$; for blocking HIGHs $\gamma=0.2$. Subtropical HIGHs are static,
centred at $(h\,\phi_H,\lambda_{\mathrm{basin}})$.
\item TROPICAL. Recurvature is a closed-form parabola in the local (east, poleward) frame at
genesis $(\phi_0,\lambda_0)$:
\begin{equation}
  x(\tau)=-u_w\tau+\tfrac12\,\alpha_r\tau^2,\quad
  y(\tau)=v_w\tau,\qquad
  \alpha_r=\frac{u_w}{0.4\,T_{\mathrm{TC}}},
\end{equation}
with $u_w=5$ and $v_w=2.5$\,m\,s$^{-1}$$^\dagger$. The storm moves west until $\tau=0.4\,T$,
then east (recurvature). It is mapped to the sphere by the exponential map at the genesis
point. Once $|\phi|>\phi_H$, or once it is over land for more than $0.1\,T$, it decays: its
strength is multiplied by $1-S(0,0.15\,T,\tau-\tau_{\mathrm{exit}})$.
\end{itemize}

\paragraph{Population (stateless renewal).}
Each hemisphere's storm track is divided into $n_s$ sectors of width $2\pi/n_s$, with
$n_s=\max(1,\operatorname{round}k^*)$. Sector $q$ has seeded slots $n\in\mathbb{Z}$:
\begin{equation}
  t_{b}(q,n)=(n+\eta_{qn})\,\Delta t,\qquad \Delta t=T_{\mathrm{LOW}},\quad
  \eta_{qn}\sim\mathcal{U}(0,1),\quad \text{spawn if } \zeta_{qn}<0.85,
\end{equation}
where $(\eta,\zeta)$ come from a hash of $(\mathrm{seed},q,n)$. Only the slots
$n\in\{\lfloor t/\Delta t\rfloor-2,\dots,\lfloor t/\Delta t\rfloor\}$ can be active, so the
active set at $t$ is computed with no state. A slot's LOW is born on the downstream side of
the trough nearest the sector centre:
\begin{equation}
  \lambda_0=\lambda_{\mathrm{trough}}(t_b)+0.25\,\frac{\lambda_E}{a\cos\phi_0},\qquad
  \phi_0 = \phi_j(\lambda_0,t_b)+h\,0.3\,L_\beta^\circ\ \text{(poleward side of the jet)}.
\end{equation}
The expected number of active LOWs per hemisphere is
$n_s\times0.85\times\bar T/\Delta t\approx0.85\,k^*$. HIGHs, TCs and CUTOFFs use the same
scheme with their own counts and slots. TC genesis points are drawn in the TC band, weighted
by the genesis proxy of \S\ref{sec:eq-rules}, over sea only.

\paragraph{Spacing.}
Objects are processed in birth order, with pinned ones first. An object is dropped if
\begin{equation}
  \operatorname{gc}(\mathbf{c}_i(t),\mathbf{c}_k(t))<0.75\,(r_i+r_k)
\end{equation}
for an object $k$ already kept, where $\operatorname{gc}$ is the great-circle distance.

%--------------------------------------------------------------------
\subsection{Pressure}
\label{sec:eq-pressure}
The climatological zonal pressure (hPa)$^\dagger$, with $\phi_{st}$ the storm-track centre, is
\begin{equation}
  p_{\mathrm{clim}}(\phi)=1013-4\,G(\phi;\phi_{\mathrm{ITCZ}},8^\circ)
  +8\,G(|\phi|;\phi_H,8^\circ)-10\,G(|\phi|;\phi_{st},10^\circ)+4\,G(|\phi|;90^\circ,10^\circ).
\end{equation}
In each object's local frame $(x,y)$, with $x$ along the jet and $y$ poleward, LOW and HIGH
kernels are anisotropic Gaussians:
\begin{equation}
  \Delta p_{\mathrm{LOW}}=-\Sigma(s)\exp\!\left[-\frac{x^2}{(1.3r)^2}-\frac{y^2}{(0.8r)^2}\right],\qquad
  \Delta p_{\mathrm{HIGH}}=+\Sigma(s)\exp\!\left[-\frac{x^2+y^2}{(1.5r)^2}\right],
\end{equation}
with $P_{\mathrm{LOW}}=25$ and $P_{\mathrm{HIGH}}=15$\,hPa$^\dagger$. Highs get broader,
weaker kernels \cite{hoskins2002}.

A tropical cyclone uses the Holland profile \cite{holland1980}:
\begin{equation}
  \Delta p_{TC}(r)=-\Sigma(s)\left(1-\exp\!\left[-\left(\frac{R_m}{r}\right)^{B}\right]\right),
  \qquad B\sim\mathcal{N}(1.3,0.35)\ \text{clamped to }[0.8,2.2]
  \ \text{\cite{willoughby2004}},
\end{equation}
with $P_{TC}=60$\,hPa$^\dagger$. The total is
$p=p_{\mathrm{clim}}+\sum_i\Delta p_i$. Each kernel is evaluated only within $4r_i$ of its
centre (spatial culling).

%--------------------------------------------------------------------
\subsection{Wind}
\label{sec:eq-wind}
The geostrophic wind, with $\rho=1.2$\,kg\,m$^{-3}$ and $f$ floored near the equator, is
\begin{equation}
  u_g=-\frac{1}{\rho\,\tilde f}\frac{\partial p}{\partial y},\qquad
  v_g=\frac{1}{\rho\,\tilde f}\frac{\partial p}{\partial x},\qquad
  \tilde f=h\,\max\!\big(|f(\phi)|,f(5^\circ)\big).
\end{equation}
Within $3R_m$ of a tropical cyclone, the gradient wind from the Holland profile replaces
it, blowing tangentially in the cyclonic sense:
\begin{equation}
  V(r)=\sqrt{\frac{B\,\Delta p}{\rho}\left(\frac{R_m}{r}\right)^{B}
  e^{-(R_m/r)^B}+\frac{r^2f^2}{4}}-\frac{rf}{2}.
\end{equation}
Friction at the surface turns the wind toward low pressure, counter-clockwise in the north,
and slows it:
\begin{equation}
  \mathbf{v}_s=0.7\,\mathsf{R}(h\cdot25^\circ)\,\mathbf{v}_g.
\end{equation}
Near the equator the wind blends into the trades:
\begin{equation}
  \mathbf{v}=w\,\mathbf{v}_s+(1-w)\,\mathbf{v}_{\mathrm{trade}},\quad
  w=S(5^\circ,12^\circ,|\phi|),\quad
  \mathbf{v}_{\mathrm{trade}}=(-5,\,-1.5\,h)\ \mathrm{m\,s^{-1}}.
\end{equation}
Derivatives are central differences on the sphere, over a step of $0.05\,r_c$.

%--------------------------------------------------------------------
\subsection{Fronts}
\label{sec:eq-fronts}
Fronts are drawn in each LOW's local frame, with $(x,y)$ in units of $r$, $x$ along its
motion and $y$ poleward. Each front is a curve of linearly varying heading, i.e.\ a circular
arc in closed form:
\begin{equation}
  \boldsymbol{\gamma}(\sigma)=\mathbf{P}_0+\int_0^{\sigma}
  \big(\cos\vartheta(\sigma'),\sin\vartheta(\sigma')\big)\,d\sigma',\qquad
  \vartheta(\sigma)=\vartheta_0+\kappa\,\sigma,\quad \sigma\in[\sigma_0,\ell].
\end{equation}
Angles are measured from $+x$ toward poleward.

\paragraph{Occlusion and triple point.} Let $o(s)=S(0.35,0.9,s)$. The occluded front runs from
the centre to the triple point $\mathbf{P}_T$, curving poleward (it wraps):
\begin{equation}
  \ell_o=1.2\,o(s),\quad \vartheta_0=-20^\circ,\quad \kappa_o=+60^\circ\!/r,\qquad
  \mathbf{P}_T=\boldsymbol{\gamma}_o(\ell_o).
\end{equation}

\paragraph{Norwegian variant} (the default):
\begin{align}
  \text{warm front}:&\quad \mathbf{P}_0=\mathbf{P}_T,\ \vartheta_0=-10^\circ,\ \kappa=-10^\circ\!/r,\
     \ell_w=1.5\,S(0,0.2,s);\\
  \text{cold front}:&\quad \mathbf{P}_0=\mathbf{P}_T,\ \vartheta_0=-80^\circ,\ \kappa=-16^\circ\!/r,\
     \ell_c=2.5\,S(0,0.2,s),\ \sigma_0=0.
\end{align}

\paragraph{Shapiro--Keyser variant} \cite{shapiro1990}, a seeded 30\,\% of LOWs:
\begin{itemize}
\item the warm front continues as a bent-back front, with
  $\vartheta(\sigma)=-10^\circ+\max(0,\sigma-1)\cdot150^\circ\!/r$;
\item the cold front is fractured, with $\sigma_0=0.4\,S(0.3,0.6,s)$;
\item the cold front meets the warm front at a T-bone, with $\vartheta_0=-90^\circ$;
\item there is no occlusion: $\ell_o=0$, and a warm seclusion $G(\rho;0,0.3)$ clears the core
  for $s>0.7$.
\end{itemize}

\paragraph{Signed distances.}
For a point $\mathbf{q}$ and front $f$, let $\sigma^*$ be the parameter of the closest point on
the front. The signed distance, positive on the cold side, is
\begin{equation}
  d_f(\mathbf{q})=\pm\min_{\sigma}|\mathbf{q}-\boldsymbol{\gamma}_f(\sigma)|.
\end{equation}
The sign is positive on the left of the warm front's heading (poleward) and on the right of
the cold front's heading (west, behind it). $\sigma^*$ also gives the along-front taper
$T_f=1-S(0.7\ell,\ell,\sigma^*)$. The warm sector is $\{d_w<0\}\cap\{d_c<0\}\cap\{\rho<2.5\}$.

%--------------------------------------------------------------------
\subsection{Genus rules}
\label{sec:eq-rules}
Each genus channel $g$ accumulates the contributions of all systems with a soft maximum,
$c_g=1-\prod_i(1-c_{g,i})$, and is clamped to $[0,1]$. Distances are in units of the system's
$r$. $\Sigma'=\Sigma(s)/P$ is the normalised intensity.

\paragraph{Warm-front shield} \cite{carlson1980}. Let $\xi=d_w$ on the cold side
($\xi\ge0$), with taper $T_w$ and edge softness $e=0.08$. Each layer is
$c_g=C_g\,\Pi(\xi;\xi_0,\xi_1,e)\,T_w\,\Sigma'$, with:

\begin{center}
\begin{tabular}{lcccl}
\toprule
genus & $\xi_0$ & $\xi_1$ & $C_g$ & channel\\
\midrule
nimbostratus & 0.00 & 0.20 & 0.95 & A2\\
altostratus  & 0.15 & 0.50 & 0.85 & A3\\
cirrostratus & 0.45 & 0.85 & 0.70 & B2\\
cirrus       & 0.80 & 1.20 & 0.55 & B1\\
\bottomrule
\end{tabular}
\end{center}

\paragraph{Cold front} \cite{browning1994}. A narrow Cb line on the front, with congestus in a
band ahead of it on the warm side:
\begin{align}
  c_{\mathrm{conv}}&=0.9\,G(d_c;0,0.08)\,T_c\,\Sigma',&
  \mathrm{depth}&=G(d_c;0,0.08)\,\big(1-0.5\,\sigma^*/\ell_c\big);\\
  c_{\mathrm{conv}}&\mathrel{+}=0.5\,\Pi(-d_c;0.08,0.25,0.05)\,T_c\,\Sigma',&
  \mathrm{depth}&\mathrel{+}=0.5\,\Pi(-d_c;0.08,0.25,0.05).
\end{align}
The anvil shelf is
$\mathrm{shelf}=\max(\mathrm{shelf},\,S(0.8,0.95,\mathrm{depth}))$, stretched downwind by the
existing \texttt{anvil\_spread}.

\paragraph{Warm sector.} Overcast with drizzle near the warm front; stratocumulus elsewhere,
mostly over sea. With $m$ the sea fraction:
\begin{equation}
  c_{St}=0.6\,G(d_w;0,0.4)\,W\,\Sigma',\qquad
  c_{Sc}=(0.35+0.3\,m)\,W\,\Sigma',\qquad
  W=\mathbf{1}_{\text{warm sector}}\cdot S(2.5,2.0,\rho).
\end{equation}

\paragraph{Comma head and dry slot} \cite{carlson1980,browning1994}.
\begin{align}
  H&=G(x;0.3,0.9)\,G(y;0.6,0.5)\,S(0.3,0.6,s)\,\Sigma':&
  c_{Ns}&\mathrel{+}=0.8H,\, c_{As}\mathrel{+}=0.6H,\, \mathrm{depth}\mathrel{+}=0.3H;\\
  D&=G(x;-0.3,0.4)\,G(y;-0.6,0.5)\,S(0.2,0.4,s):&
  c_g&\leftarrow c_g\,(1-0.9D).
\end{align}
The dry-slot factor $D$ multiplies every genus channel.

\paragraph{Post-frontal cold air.} Open-cell convection over sea, behind the cold front:
\begin{equation}
  Q=\Pi(d_c;0.2,2.0,0.2)\,m\,\Sigma':\qquad
  c_{\mathrm{conv}}\mathrel{+}=0.45Q,\ \mathrm{depth}\mathrel{+}=0.3Q,\ c_{Sc}\mathrel{+}=0.3Q.
\end{equation}

\paragraph{Anticyclones} \cite{wood2012,klein1993}. Let $\rho$ be the distance from the centre in
units of $r$, $T$ the climate temperature (\textdegree C) and $m$ the sea fraction.
\begin{align}
  \text{subsidence clearing:}\quad & c_g\leftarrow c_g\,\big(1-0.9\,\Sigma'\,G(\rho;0,0.6)\big)
     \quad g\in\{\mathrm{conv},St,Sc,Ns,As,Ac\};\\
  \text{stability proxy:}\quad & \Lambda=\Sigma'\,\big(1-S(14,26,T)\big)\,m\,
     \Big(\tfrac12+\tfrac12 S(-0.2,0.6,x)\Big);\\
  \text{Sc deck:}\quad & c_{Sc}\mathrel{+}=0.85\,\Lambda\,\Pi(\rho;0.5,1.6,0.3);\\
  \text{St / fog near the coast:}\quad & c_{St}\mathrel{+}=0.7\,\Lambda\,
     S(-200,-20,e_{\mathrm{elev}})\,\Pi(\rho;0.5,1.6,0.3);\\
  \text{flank Ac:}\quad & c_{Ac}\mathrel{+}=0.35\,\Sigma'\,\Pi(\rho;1.0,1.6,0.2).
\end{align}
The factor $\frac12+\frac12S(\dots,x)$ biases the deck to the eastern flank. In the stability
proxy, colder sea means stronger stability, about $+6\,\%$ cover per K \cite{klein1993}.

\paragraph{Tropical cyclone.} With $r$ the distance in metres, $\vartheta$ the azimuth and
$\alpha_b=15^\circ$ the band pitch angle:
\begin{align}
  \text{eye:}\quad & E=S(0.3R_m,0.6R_m,r),\qquad c_g\leftarrow E\,c_g\ \ \forall g;\\
  \text{eyewall:}\quad & c_{\mathrm{conv}}\mathrel{+}=0.95\,G(r;R_m,0.5R_m)\,\Sigma',\quad
     \mathrm{depth}\mathrel{+}=G(r;R_m,0.5R_m);\\
  \text{spiral bands:}\quad & B=\Big[\tfrac12+\tfrac12\cos\!\Big(2\vartheta-
     \frac{2\,h\ln(r/R_m)}{\tan\alpha_b}\Big)\Big]^3G(r;0,4R_m)\,S(1.5R_m,2R_m,r),\\
  & c_{\mathrm{conv}}\mathrel{+}=0.6B\,\Sigma',\quad \mathrm{depth}\mathrel{+}=0.7B;\\
  \text{CDO and outflow:}\quad & c_{Cs}\mathrel{+}=0.8\,G(r;0,3R_m)\,\Sigma'.
\end{align}

\paragraph{Air-mass convection, the ITCZ and the genesis proxy.} With humidity $H\in[0,1]$:
\begin{align}
  I&=S(0.4,0.8,H)\,S(20,30,T),\\
  c_{\mathrm{conv}}&\mathrel{+}=0.4\,I\,\big(0.5+0.5\,G(\phi;\phi_{\mathrm{ITCZ}},6^\circ)\big),\qquad
  \mathrm{depth}=\max\big(\mathrm{depth},\,I\big),\\
  \mathrm{GP}&=\big|\sin\phi\big|^{3/2}\,H^3\,S(24,28,T)\,m.
\end{align}
GP is a stand-in for the Camargo et al.\ index \cite{camargo2007}: $|\eta|\propto|\sin\phi|$,
there is no shear term, and sea temperature is replaced by the climate temperature.

\paragraph{Climate terms (kept from today's \texttt{weather.py}).} Let $W(H)$ be the existing
wetness smoothstep and $\mathbf{v}_s$ the surface wind of \S\ref{sec:eq-wind}:
\begin{align}
  F_{\mathrm{clim}}&=\big(d_f+(1-d_f)W(H)\big)\big(1-D_{\mathrm{des}}\,S_{\mathrm{hot}}\,S_{\mathrm{dry}}\big),\\
  c_{\mathrm{oro}}&=O\tanh\!\left(\frac{e(\mathbf{q})-e(\mathbf{q}-\Delta\,\hat{\mathbf{v}}_s)}
     {\Delta\,s_0}\right)W(H),\\
  c_{\mathrm{conv}}&=\operatorname{clamp}\Big(Z(\phi)F_{\mathrm{clim}}+c_{\mathrm{oro}}+M\,m
     +A\,\operatorname{fbm}(\mathbf{q}_{\mathrm{warped}})+\textstyle\sum c_{\mathrm{conv},i},\,0,1\Big)
     \cdot E\cdot(1-\text{clearings}).
\end{align}
The background fair-weather cover follows the belts:
\begin{equation}
  Z(\phi)=0.15+0.5\,G(\phi;\phi_{\mathrm{ITCZ}},\phi_H/3)+0.35\,G(|\phi|;\phi_{st},12^\circ).
\end{equation}

\paragraph{Coverage to fbm shift (unchanged).} For each genus table,
\begin{equation}
  \operatorname{shift}_g(c)=\theta_{lo}-Q_{\mathrm{peaks}}(1-c),
\end{equation}
where $Q$ is the quantile of the sorted column peaks. This keeps the near field and the far
deck consistent.

%--------------------------------------------------------------------
\subsection{Local presets}
\label{sec:eq-pins}
A pin is $(\phi^*,\lambda^*,\text{preset},\vartheta_h)$, all at time $t$. Each preset fixes an
object kind, a target stage $s^*$ and a role point $\mathbf{X}$ in the object's local frame
(units of $r$):

\begin{center}
\begin{tabular}{llcp{6.3cm}}
\toprule
preset & object & $s^*$ & role point $\mathbf{X}$\\
\midrule
FAIR\_HIGH & HIGH & 0.5 & $(0,0)$\\
WARM\_FRONT\_APPROACH & LOW & 0.4 & $\boldsymbol\gamma_w(0.5\ell_w)+0.3\,\mathbf{n}_w$ (As overhead)\\
COLD\_FRONT & LOW & 0.4 & $\boldsymbol\gamma_c(0.4\ell_c)$\\
SQUALL\_LINE & LOW & 0.4 & $\boldsymbol\gamma_c(0.4\ell_c)-0.17\,\mathbf{n}_c$ (congestus/Cb band)\\
OVERCAST\_DRIZZLE & LOW & 0.4 & $\boldsymbol\gamma_w(0.3\ell_w)-0.2\,\mathbf{n}_w$ (warm sector)\\
HURRICANE & TROPICAL & 0.5 & $(0,0)$\\
THUNDERSTORMS & none & --- & no object: $I\leftarrow\max(I,0.9)$ within $0.3\,r_c$\\
\bottomrule
\end{tabular}
\end{center}

Here $\mathbf{n}_f$ is the unit normal toward the front's cold side.

\paragraph{Placing the pinned object.} With $\mathsf{R}(\vartheta_h)$ the rotation by the
heading (default 0), the centre is the exponential map at the pin:
\begin{equation}
  \mathbf{c}=\exp_{\mathbf{p}^*}\!\big(-r\,\mathsf{R}(\vartheta_h)\,\mathbf{X}\big).
\end{equation}
The track is then inverted so that the object is at $\mathbf{c}$ with stage $s^*$ at time $t$:
\begin{equation}
  t_b=t-s^*T,\qquad
  \lambda_0=\lambda_c-\frac{\gamma c_s\,s^*T}{a\cos\phi_0},\qquad
  \phi_0=\phi_c-h\frac{v_p\,s^*T}{a}.
\end{equation}
For a TROPICAL pin, the parabola is inverted for the genesis point. If that point is over
land or outside the band, the pin is still accepted: the storm is then taken as already
mature at the pin.

\paragraph{Refitting the jet phases.} Write each wave as
$a_m\sin k_m\lambda'+b_m\cos k_m\lambda'$, with $\lambda'$ including $-\varpi_m t$. A pinned
LOW asks for a trough upstream of it,
$\lambda_{\mathrm{tr}}=\lambda_c-0.25\lambda_E/(a\cos\phi_c)$, with target
$\phi_j(\lambda_{\mathrm{tr}})=\phi_{0,j}-hA_{\mathrm{tot}}$. A pinned HIGH asks for a ridge,
$\phi_j(\lambda_c)=\phi_{0,j}+hA_{\mathrm{tot}}$. The coefficients solve the ridge regression
\begin{equation}
  \min_{a,b}\ \sum_{i\in\text{pins}}\big(\phi_j(\lambda_i)-\phi_i^{\mathrm{target}}\big)^2
  +\mu\sum_m\big[(a_m-a_m^0)^2+(b_m-b_m^0)^2\big],\qquad \mu=0.1,
\end{equation}
which is linear in $(a,b)$, so it has a closed-form solution. The result is rescaled so that
\begin{equation}
  \sum_m\sqrt{a_m^2+b_m^2}\le1.5\,A_{\mathrm{tot}}.
\end{equation}

\paragraph{Feasibility (raise, never move).}
\begin{align}
  \text{HURRICANE:}\quad & m(\mathbf{p}^*)>0.5\ \wedge\ |\phi^*|\in\text{TC band}\pm2^\circ
     \ \wedge\ T(\mathbf{p}^*)\ge24^\circ\mathrm{C};\\
  \text{LOW presets:}\quad & |\phi^*|\in[\phi_{st,lo}-L_R^\circ,\ \phi_{st,hi}+L_R^\circ];\\
  \text{THUNDERSTORMS:}\quad & S(0.4,0.8,H)\,S(20,30,T)\ge0.3;\\
  \text{two pins:}\quad & \operatorname{gc}(\mathbf{c}_i,\mathbf{c}_k)\ge0.75(r_i+r_k),
\end{align}
unless both pins ask for the same kind and their centres fall within $0.5\,r$, in which case
they share one object (for example COLD\_FRONT and SQUALL\_LINE on one low).

%====================================================================
\section{Open points}
\begin{itemize}
\item Warm-front distances and slopes vary by source (1:100--1:300). We tune them on preview
  maps within that range.
\item The GP-like genesis proxy uses the terrain climate fields, not real SST or shear.
\item The content of \cite{hillaire2016,bauer2019,harris2003,browning1990} was not checked
  beyond the bibliographic data.
\end{itemize}

\bibliographystyle{plainnat}
\bibliography{synoptic-weather}
\end{document}
